% papers/torre_induc.tex — prova de indução dos tecidos (T+T*) na espiral.
% Escalada formal: papers/corpo_peano.tex
\documentclass[11pt,a4paper]{article}
\usepackage[utf8]{inputenc}\usepackage[T1]{fontenc}\usepackage[portuguese]{babel}
\usepackage{amsmath,amssymb}
\usepackage[margin=2.5cm]{geometry}
\newcommand{\code}[1]{\texttt{#1}}
\begin{document}
\noindent\textbf{Torre por indução} --- o passo $T_{k+1}=T_k+T_k^{*}$ sobe um andar de cada vez;
cada \texttt{\^{}\{...\}} é um giro do inversor.
\[
  a^{\,b^{\,c^{\,d^{\,e^{\,f^{\,g}}}}}}\,.
\]
Seis níveis: dimensão da torre $d_6=128$, interface $\iota_6=24$
(família $6\to12\to24$). A dimensão da torre \textbf{não} é o índice da lei.

\medskip\noindent Quatro objectos, relacionados e não identificados
(\code{papers/corpo\_peano.tex}):
\[
\underbrace{\mathbb{Z}/8\mathbb{Z}}_{\text{leis}}
\qquad
\underbrace{(\mathbb{N}_{>0},\mid)}_{\text{períodos}}
\qquad
\underbrace{d_k=2^{k+1}}_{\text{torre}}
\qquad
\underbrace{C_k=12+6\lfloor k/3\rfloor}_{\text{régua}}.
\]
O período oito é do catálogo ($\ell_{m+8}=\ell_m$); a torre cresce.
Verificação: \code{tests/corpo\_peano.c}, \code{tests/torre\_induc.js}.
\end{document}
